\section*{Capítulo \ref{ch:g11:lenses-and-eye} --- Lentes, imagens e o olho}

\begin{solution}{exo:g11:lenses-and-eye:1}
$C = 1/0.50 = +2.0$~D; $1/(-0.20) = -5.0$~D; $1/0.0040 = +250$~D. Uma
\omterm{def:g11:lenses-and-eye:lens}{lente} de $+8.0$~D: $f' = 1/8.0 = \qty{0.125}{m} = \qty{12.5}{cm}$.
\end{solution}

\begin{solution}{exo:g11:lenses-and-eye:2}
\omterm{def:g11:lenses-and-eye:lens}{Divergente} ($f' < 0$). Não: para qualquer distância $d$ do objeto,
$\gamma = f'/(f' - d)$ (mostrado no \cref{exo:g11:lenses-and-eye:15})
fica entre $0$ e $1$ --- a imagem é sempre \omterm{def:g11:lenses-and-eye:images}{virtual}, direita e
\emph{reduzida}. A letra impressa aparece direita e encolhida, nunca
ampliada.
\end{solution}

\begin{solution}{exo:g11:lenses-and-eye:3}
O objeto está em $2f'$; os raios central e paralelo se cruzam
\qty{20}{cm} atrás da \omterm{def:g11:lenses-and-eye:lens}{lente}. Conferindo:
$1/\overline{OA'} = 1/10 - 1/20 = 1/20$, logo
$\overline{OA'} = +\qty{20}{cm}$ e $\gamma = 20/(-20) = -1$: \omterm{def:g11:lenses-and-eye:images}{real},
invertida e do mesmo tamanho --- como no desenho.
\end{solution}

\begin{solution}{exo:g11:lenses-and-eye:4}
$1/\overline{OA'} = 1/8.0 - 1/24 = 1/12$:
$\overline{OA'} = +\qty{12}{cm}$, $\gamma = 12/(-24) = -0.50$. \omterm{def:g11:lenses-and-eye:images}{Imagem
real} \qty{12}{cm} atrás da \omterm{def:g11:lenses-and-eye:lens}{lente}, invertida, com metade do tamanho.
\end{solution}

\begin{solution}{exo:g11:lenses-and-eye:5}
A \omterm{def:g11:lenses-and-eye:eye}{córnea} e o \omterm{def:g11:lenses-and-eye:eye}{cristalino} formam a \omterm{def:g11:lenses-and-eye:lens}{lente convergente}; a \omterm{def:g11:lenses-and-eye:eye}{retina} é o
anteparo. A \omterm{def:g11:lenses-and-eye:eye}{acomodação} muda a \omterm{def:g11:lenses-and-eye:focal}{distância focal} (a \omterm{def:g11:lenses-and-eye:vergence}{vergência}) do
\omterm{def:g11:lenses-and-eye:eye}{cristalino}; a distância lente--retina não pode mudar. \omterm{def:g11:lenses-and-eye:eye}{Ponto próximo}: o
ponto mais próximo visto com nitidez com \omterm{def:g11:lenses-and-eye:eye}{acomodação} máxima ---
\qty{25}{cm} para o \omterm{def:g11:lenses-and-eye:eye}{olho padrão}.
\end{solution}

\begin{solution}{exo:g11:lenses-and-eye:6}
Paisagem: sensor no plano focal, a \qty{50}{mm} da \omterm{def:g11:lenses-and-eye:lens}{lente}. A \qty{3.0}{m}:
$1/\overline{OA'} = 1/50 - 1/3000 = 59/3000$, logo
$\overline{OA'} \approx \qty{50.8}{mm}$. Curso: cerca de \qty{0.8}{mm}.
\end{solution}

\begin{solution}{exo:g11:lenses-and-eye:7}
$1/\overline{OA'} = 1/6.0 - 1/5.0 = -1/30$:
$\overline{OA'} = \qty{-30}{cm}$ --- \omterm{def:g11:lenses-and-eye:images}{virtual}, direita,
$\gamma = (-30)/(-5.0) = +6.0$: seis vezes o selo, \qty{30}{cm} acima da
\omterm{def:g11:lenses-and-eye:lens}{lente}.
\end{solution}

\begin{solution}{exo:g11:lenses-and-eye:8}
$1/f' = 1/40 + 1/40 = 1/20$: $f' = \qty{20}{cm}$;
$\gamma = 40/(-40) = -1$. A simetria obriga
$\abs{\overline{OA'}} = \abs{\overline{OA}}$ e, portanto,
$\abs{\gamma} = 1$; a imagem é \omterm{def:g11:lenses-and-eye:images}{real} e invertida, logo $\gamma = -1$:
exatamente do mesmo tamanho.
\end{solution}

\begin{solution}{exo:g11:lenses-and-eye:9}
$1/\overline{OA'} = -1/25 - 1/50 = -3/50$:
$\overline{OA'} \approx \qty{-16.7}{cm}$,
$\gamma = (-16.7)/(-50) = +0.33$: \omterm{def:g11:lenses-and-eye:images}{virtual}, direita, com um terço do
tamanho. Os óculos de um \omterm{def:g11:lenses-and-eye:defects}{míope} são \omterm{def:g11:lenses-and-eye:lens}{divergentes}: tudo o que se vê através
deles --- inclusive, de fora, os olhos de quem os usa --- fica reduzido.
\end{solution}

\begin{solution}{exo:g11:lenses-and-eye:10}
Imagem do infinito no \omterm{def:g11:lenses-and-eye:eye}{ponto remoto}: $f' = \qty{-0.40}{m}$, logo
$C = -2.5$~D. A imagem da montanha fica \qty{40}{cm} à frente da \omterm{def:g11:lenses-and-eye:lens}{lente}
--- exatamente no \omterm{def:g11:lenses-and-eye:eye}{ponto remoto}, que o olho relaxado vê com nitidez e sem
\omterm{def:g11:lenses-and-eye:eye}{acomodação} alguma.
\end{solution}

\begin{solution}{exo:g11:lenses-and-eye:11}
$C = 1/\overline{OA'} - 1/\overline{OA} = 1/(-0.80) - 1/(-0.25)
= -1.25 + 4.00 = +2.75$~D.
\end{solution}

\begin{solution}{exo:g11:lenses-and-eye:12}
\emph{1.} $\gamma = \overline{OA'}/\overline{OA} = -50$ dá
$\overline{OA'} = -50\,\overline{OA}$; com
$\overline{OA'} - \overline{OA} = \qty{5.1}{m}$,
$-51\,\overline{OA} = \qty{5.1}{m}$: $\overline{OA} = \qty{-0.10}{m}$ e
$\overline{OA'} = +\qty{5.0}{m}$.

\emph{2.} $1/f' = 1/5.0 + 1/0.10 = 10.2$: $f' \approx \qty{9.8}{cm}$.

\emph{3.} $50 \times \qty{24}{mm} = \qty{1.2}{m}$, invertida --- e é por
isso que os slides são colocados de cabeça para baixo.
\end{solution}

\begin{solution}{exo:g11:lenses-and-eye:13}
\emph{1.} $C_\infty = 1/0.017 \approx 59$~D;
$C_{25} = 1/0.017 + 1/0.25 \approx 63$~D. \omterm{def:g10:signals-and-waves:amplitude}{Amplitude}:
$1/0.25 = 4.0$~D.

\emph{2.} \omterm{def:g10:signals-and-waves:amplitude}{Amplitude} $= 1/1.0 = 1.0$~D: resta um quarto ($25\%$) dos
$4.0$~D do \omterm{def:g11:lenses-and-eye:eye}{olho padrão}.
\end{solution}

\begin{solution}{exo:g11:lenses-and-eye:14}
\emph{1.} $f' = 0.25/3 \approx \qty{0.083}{m} = \qty{8.3}{cm}$.

\emph{2.} $\overline{OA'} = \qty{-25}{cm}$:
$1/\overline{OA} = -1/25 - 3/25 = -4/25$, logo
$\overline{OA} = \qty{-6.25}{cm}$ e $\gamma = (-25)/(-6.25) = +4.0$ ---
melhor do que o $\times 3$ anunciado, que supõe a imagem no infinito;
puxá-la até o \omterm{def:g11:lenses-and-eye:eye}{ponto próximo} ganha uma \omterm{def:g10:orders-of-magnitude:unit}{unidade}.
\end{solution}

\begin{solution}{exo:g11:lenses-and-eye:15}
$1/\overline{OA'} = 1/f' - 1/d = (d - f')/(f'd)$, logo
$\overline{OA'} = f'd/(d - f')$ e
$\gamma = \overline{OA'}/(-d) = f'/(f' - d)$. Se $d > f'$:
$\overline{OA'} > 0$ (\omterm{def:g11:lenses-and-eye:images}{real}) e $\gamma < 0$ (invertida). Se
$0 < d < f'$: $\overline{OA'} < 0$ (\omterm{def:g11:lenses-and-eye:images}{virtual}) e
$\gamma = f'/(f' - d) > 1$: direita e ampliada --- a configuração de
\omterm{def:g11:lenses-and-eye:magnifier}{lupa}.
\end{solution}

\begin{solution}{pb:g11:lenses-and-eye:1}
\textbf{1.} $1/\overline{OA'} = 1/20 - 1/60 = 1/30$:
$\overline{OA'} = +\qty{30}{cm}$, $\gamma = -0.50$ --- \omterm{def:g11:lenses-and-eye:images}{real}, invertida,
com metade do tamanho.

\textbf{2.} $\overline{OA'} = +\qty{60}{cm}$, $\gamma = -2$: as posições
\qty{30}{cm} e \qty{60}{cm} trocaram de papel e os \omterm{def:g11:lenses-and-eye:magnification}{aumentos} são inversos
um do outro --- a luz invertida percorre o mesmo caminho.

\textbf{3.} $1/\overline{OA'} = 1/20 - 1/10 = -1/20$:
$\overline{OA'} = \qty{-20}{cm}$, $\gamma = +2$ --- \omterm{def:g11:lenses-and-eye:images}{virtual}, direita,
duplicada.

\textbf{4.} Partindo do infinito, a imagem começa no plano focal; à
medida que o objeto se aproxima de $F$, a \omterm{def:g11:lenses-and-eye:images}{imagem real} e invertida se
afasta rumo ao infinito e cresce; dentro de $f'$ ela se torna \omterm{def:g11:lenses-and-eye:images}{virtual},
direita e ampliada, à frente da \omterm{def:g11:lenses-and-eye:lens}{lente}.

\textbf{5.} $1/\overline{OA} \approx 0$ deixa
$1/\overline{OA'} = 1/f'$: $\overline{OA'} = f'$, o plano focal.

\textbf{6.} No plano focal, \qty{50}{mm} atrás da \omterm{def:g11:lenses-and-eye:lens}{lente}.

\textbf{7.} $1/\overline{OA'} = 1/50 - 1/1000 = 19/1000$:
$\overline{OA'} \approx \qty{52.6}{mm}$; curso
$\approx \qty{2.6}{mm}$.

\textbf{8.} $1/\overline{OA} = 1/55 - 1/50 = -1/550$: distância mínima de
focalização \qty{550}{mm} = \qty{55}{cm}.

\textbf{9.} $\gamma = 55/(-550) = -0.10$: o rosto tem imagem de
\qty{20}{mm}, cabendo com folga no sensor de \qty{24}{mm}.

\textbf{10.} $\gamma = -1$ obriga
$\overline{OA'} = -\overline{OA}$, logo $2/\overline{OA'} = 1/f'$:
objeto a \qty{100}{mm} ($= 2f'$) à frente e sensor a \qty{100}{mm} atrás
--- o dobro do ajuste de paisagem, muito além do mecanismo de
\qty{55}{mm}; as objetivas macro fornecem o curso extra.

\textbf{11.} Tudo o que estiver mais perto do que o \omterm{def:g11:lenses-and-eye:eye}{ponto próximo} dela
(\qty{1.0}{m}) fica borrado: os braços esticados aproximam a página de
\qty{1.0}{m} o máximo possível --- nítida, ainda que pequena.

\textbf{12.} $C = 1/(-1.0) - 1/(-0.25) = -1.0 + 4.0 = +3.0$~D.

\textbf{13.} $\gamma = (-1.0)/(-0.25) = +4.0$: quatro vezes maior, mas
quatro vezes mais longe --- o mesmo tamanho angular; o ganho não é de
tamanho, e sim de \emph{nitidez}, já que a imagem agora está no \omterm{def:g11:lenses-and-eye:eye}{ponto
próximo} dela.

\textbf{14.} Para ficar nítida, a imagem precisa estar além de
\qty{1.0}{m}: da página a \qty{25}{cm} (imagem a \qty{1.0}{m}) até a
página a $f' = 1/3.0 \approx \qty{0.33}{m}$ (imagem no infinito). Um
objeto distante tem imagem \qty{33}{cm} \emph{atrás} dos óculos --- uma
\omterm{def:g11:lenses-and-eye:images}{imagem real} que o olho não consegue usar: o mundo distante fica borrado.

\textbf{15.} Daí as \omterm{def:g11:lenses-and-eye:lens}{lentes} bifocais: a correção de leitura na metade de
baixo da \omterm{def:g11:lenses-and-eye:lens}{lente} e nenhuma correção (ou a de longe) na metade de cima.

\textbf{16.} $1/\overline{OA'} = 1/5.0 - 1/4.0 = -1/20$:
$\overline{OA'} = \qty{-20}{cm}$, \omterm{def:g11:lenses-and-eye:images}{virtual}, direita, $\gamma = +5.0$.

\textbf{17.} $1/\overline{OA} = -1/25 - 1/5.0 = -6/25$:
$\overline{OA} \approx \qty{-4.2}{cm}$;
$\gamma = (-25)/(-25/6) = +6.0$.

\textbf{18.} Pedra em $F$: imagem no infinito, observada com \omterm{def:g11:lenses-and-eye:eye}{acomodação}
nula --- horas de exame sem cansaço visual.

\textbf{19.} $C_\infty = 1/0.017 \approx 59$~D;
$C_{25} = 1/0.017 + 1/0.25 \approx 63$~D: cerca de $4$~D de \omterm{def:g11:lenses-and-eye:eye}{acomodação}.

\textbf{20.} $f' = 1/58.8 \approx \qty{17.0}{mm}$ contra
$1/62.8 \approx \qty{15.9}{mm}$: cerca de \qty{1.1}{mm}, uns $6\%$ --- o
zoom embutido que a câmera imitou com \qty{2.6}{mm} de curso e que a avó
remendou com $+3$~D.
\end{solution}
